\scanpage{269}
en appliquant $\mathbb R\varphi_{\unsure{0|N_0},\infty}$

\[
\underbrace{(\mathbb R^*\varphi_{N_0,\infty})(\tau_{\infty,\deleted{\unsure{i}}N_0}i)(K)}
\simeq(\mathbb R^{*}\varphi_{N_0,\infty})(\tau_{\infty N_0}L^\cdot)
\]

\[
\text{i.e.}\qquad i(K)\simeq
\underbrace{(\mathbb R^*\varphi_{N_0,\infty})(\tau_{\infty,N_0}L^\cdot)}_{L'^\cdot}
\]

\deleted{Je dis que l'hom d'augmentation}

Mais c'est là un isom dans $D(\mathcal M_{N_0})$ \unsure{seult|seulement}, -- rien de plus précis que la donnée de départ -- on aimerait un isom dans $D(C)$. On va supposer $i$ pl.\ fidèle et avec un adjoint à droite $k$,

\[
\text{donc}\qquad k\circ i\simeq\mathrm{id}_C
\]

et on va supposer $k$ \uline{exact}. \unsure{Appliquant|Appliquons} $k$ à l'isom obtenu, on trouve

\[
K^\cdot\simeq k\bigl(\mathbb R^*\varphi_{N_0,\infty}(\tau_{\infty,N_0}L)\bigr)
\]

\inserted{$\simeq(\mathbb R^*(k\varphi_{N_0\infty}))(\tau_{\infty N_0}L^\cdot)$}

i.e. $K^\cdot$ se \unsure{reconstitue} dans $D(\unsure{\mathcal C|C})$ à l'aide de $\tau_{\infty,N_0}(i(K^\cdot))$ \deleted{\unsure{dans}} $\tau_{\infty,N_0}(L^\cdot)$ dans $D(\mathcal M_\infty)$.

\deleted{On va simplifier l'expression obtenue, en considérant}

\[
\deleted{\varphi_{N_0,\infty}(\tau_{\infty,N_0}(L^\cdot))\longrightarrow\mathbb R^*\varphi_{N_0,\infty}(\tau_{\infty N_0}L))}
\]

\deleted{\uline{Lemme} Cet hom est un quasi-isomorphisme}

\deleted{\uline{\unsure{Corollaire}} $\begin{aligned} K^\cdot&\simeq k\varphi_{N_0,\infty}(\tau_{\infty,N_0}L^\cdot)\\ &\simeq k\varphi_{N_0,\infty}(\tau_{\infty N_0}iK).\end{aligned}$}

% archive sheet number: 268
